\newcommand{\ModuleID}{EL-M14}
\newcommand{\ModuleIDKey}{M14}
\newcommand{\ModuleTitle}{Orations and Petitions}
\newcommand{\ModuleStage}{Missal Reading}
\newcommand{\ModuleUnit}{III.3,III.4}
\newcommand{\ModuleOutcome}{Unfold the compressed architecture of Collects, Secrets, and Postcommunions and compare it with longer periodic prose}
\newcommand{\ModulePrerequisites}{Independent reader workflow and Ordinary navigation through EL-M13}
\newcommand{\ModuleNext}{EL-M15, Lessons, Epistles, and Gospels}
\newcommand{\ModuleBeforeContinuing}{Map address, ground, petition, offering or sacramental context, intended effect, and conclusion without forcing ambiguous clause roles; rebuild an oration from its syntax; compare its compression with Cicero or Augustine; and compose a clearly non-liturgical civic petition from original propositions with documented frames.}
\newcommand{\ModuleDelayNote}{}
\newcommand{\ModuleExtensionPractice}{Return to the exact memory set, reader passage, integrated practice, or worksheet that exposes the uncertainty; then use different Latin for the later return.}
\newcommand{\ModuleLessonSource}{\CourseRoot/shared/content/volume3/all}
\newcommand{\ModuleExerciseSource}{\CourseRoot/shared/exercises/volume3_4/volume3}
\newcommand{\ModuleWorkedBank}{III}
\newcommand{\ModuleExerciseBank}{III}
\newcommand{\ModuleWorksheetVariants}{A,B,C,D}
\newcommand{\ModuleMemorySet}{M14}
\newcommand{\ModuleReaderSet}{M14}
\newcommand{\ModuleScopeNote}{Project composition is explicitly instructional and is never presented as received liturgical prayer.}
\newcommand{\ModulePractice}{%
  \begingroup
  \SelectCourseUnits{M14}
  % Integrated practice for the substantial packet sequence.  Each block joins
% the packet's underlying teaching units and is rendered twice by the packet
% shell: first as work, later as models.  Ordinary uncited Latin here is
% project-composed instructional text (P).

\begin{courseunit}{M01}
\SubmoduleExtension
  {From form to function}
  {Case endings and agreement make it possible to recover sentence roles even
  when Latin order differs from English order.}
  {Make a form--function map for \Latin{vulpes fusca parvum murum videt}.
  Reorder the Latin twice without changing its proposition, and explain why
  \Latin{vulpes} rather than \Latin{murum} is the subject.}
  {From memory, name the six ordinary cases and give one central use of each.}
  {\Latin{Vulpes} is nominative singular and controls \Latin{videt};
  \Latin{fusca} agrees with it; \Latin{murum} is accusative singular and
  \Latin{parvum} agrees with that object.  Any order preserving those forms is
  grammatically defensible, though emphasis may change.  The six cases are
  nominative, genitive, dative, accusative, ablative, and vocative.}
\SubmoduleExtension
  {Dictionary forms and visible evidence}
  {A citation form is compressed morphological information, not merely an
  English label.}
  {Expand \Latin{vulpes, vulpis f.}, \Latin{fuscus, -a, -um}, and
  \Latin{video, videre, vidi, visum}.  State which listed form reveals each
  stem and then parse every word in \Latin{vulpem fuscam video}.}
  {Explain the difference between a word's form and its work in a clause.}
  {\Latin{vulpis} reveals stem \Latin{vulp-}; the adjective supplies three
  gender forms; the infinitive supplies the present stem and conjugation,
  \Latin{vidi} the perfect stem, and \Latin{visum} the participial base.  In
  the clause, \Latin{vulpem fuscam} is accusative singular feminine object and
  \Latin{video} is first-person singular present active indicative.}
\SubmoduleExtension
  {First composition ladder}
  {Changing one feature at a time makes agreement and case choices visible.}
  {Begin with \English{the fox is brown}.  Write Latin for it; then make the
  fox plural; then replace the predicate adjective with a transitive action,
  \English{the brown foxes see the small wall}.  Label every changed ending.}
  {Cover the work and reconstruct the final sentence once.}
  {Models: \Latin{vulpes fusca est}; \Latin{vulpes fuscae sunt};
  \Latin{vulpes fuscae murum parvum vident}.  In the last sentence the subject
  and its adjective are nominative plural feminine, the object and its
  adjective accusative singular masculine, and the verb third-person plural.}
\end{courseunit}

\begin{courseunit}{M02}
\SubmoduleExtension
  {Declension as a recognition system}
  {First-, second-, and third-declension nouns must be recognized from their
  full citation forms rather than guessed from a nominative ending.}
  {For \Latin{via, viae f.}, \Latin{hortus, horti m.}, \Latin{donum, doni n.},
  \Latin{rex, regis m.}, and \Latin{animal, animalis n.}, give stem,
  declension, gender, and complete singular paradigm.}
  {State both neuter rules without consulting a table.}
  {The stems are \Latin{vi-}, \Latin{hort-}, \Latin{don-}, \Latin{reg-}, and
  \Latin{animal-}.  Neuter nominative and accusative match, and neuter
  nominative/accusative plural ends in \Latin{-a} (or the relevant third-
  declension neuter plural pattern).  Full paradigms should agree with the
  packet tables.}
\SubmoduleExtension
  {Syncretism resolved in context}
  {One visible ending can represent several case-and-number possibilities;
  syntax, agreement, and government decide among them.}
  {List every morphological possibility for \Latin{regi}, \Latin{animalia},
  and \Latin{viae}.  Use each form in one short project-composed clause that
  makes a single analysis clear.}
  {Name the evidence that rules out each unused alternative.}
  {Examples include \Latin{donum regi datur} (dative singular),
  \Latin{animalia in silva currunt} (nominative plural neuter), and
  \Latin{finis viae apparet} (genitive singular).  Other clauses are sound
  when verb agreement, government, or phrase attachment makes the analysis
  explicit.}
\SubmoduleExtension
  {Connected ordinary prose}
  {Paradigms become reading knowledge when forms are recovered inside a
  connected scene.}
  {Parse and translate: \Latin{Agricola per viam ad oppidum ambulat. In horto
  pueri animal parvum vident. Rex oppidi pueris dona dat.}  Then change the
  first sentence to plural and the second to the perfect tense.}
  {Reconstruct the three original sentences from an English proposition map.}
  {Translation: the farmer walks along the road to the town; in the garden the
  boys see a small animal; the king of the town gives gifts to the boys.
  Transformations: \Latin{agricolae per vias ad oppidum ambulant};
  \Latin{in horto pueri animal parvum viderunt}.}
\end{courseunit}

\begin{courseunit}{M03}
\SubmoduleExtension
  {The remaining noun systems}
  {Fourth and fifth declensions complete the ordinary noun map and prevent
  familiar endings from being assigned to the wrong paradigm.}
  {Decline \Latin{manus, manus f.}, \Latin{cornu, cornus n.}, \Latin{dies,
  diei m.}, and \Latin{res, rei f.}; circle forms that are morphologically
  ambiguous and resolve each in a short clause.}
  {State how citation form and gender distinguish \Latin{manus} from a
  second-declension guess.}
  {The genitives \Latin{manus} and \Latin{diei/rei}, together with gender,
  identify the systems.  A complete response reproduces the packet paradigms
  and uses agreement or syntax to resolve every circled form.}
\SubmoduleExtension
  {Agreement across unlike endings}
  {A noun and adjective agree in grammatical features, not necessarily in
  visible letters or declension.}
  {Make \Latin{fortis, forte}, \Latin{magnus, -a, -um}, and \Latin{celer,
  celeris, celere} agree with six supplied noun forms: \Latin{manui, rerum,
  cornua, regem, animalibus, die}.  Give every parse.}
  {Explain why matching endings is not a safe agreement rule.}
  {Models include \Latin{manui forti}, \Latin{rerum magnarum}, \Latin{cornua
  magna}, \Latin{regem celerem}, \Latin{animalibus fortibus}, and
  \Latin{die celere}.  Case, number, and gender match even when the endings do
  not look alike.}
\SubmoduleExtension
  {Description becoming narration}
  {Predicate adjectives, attributive adjectives, and adverbs serve different
  jobs in connected prose.}
  {Translate and transform: \Latin{Dies serenus est. Vulpes celeris per
  agrum lente ambulat, sed canis magnus eam videt.}  Make the day stormy, the
  fox plural, and the dog the object of a new sentence.}
  {Label every adjective's agreement controller and every adverb's verb.}
  {A sound version may read \Latin{dies tempestuosus est; vulpes celeres per
  agrum lente ambulant; vulpes canem magnum vident}.  \Latin{serenus} is a
  predicate adjective, \Latin{celeris/magnus} attributive, and \Latin{lente}
  modifies \Latin{ambulat}.}
\end{courseunit}

\begin{courseunit}{M04}
\SubmoduleExtension
  {Being, ability, and ordinary action}
  {\Latin{sum} and \Latin{possum} organize complements and infinitives, while
  the four regular conjugations organize present-system action.}
  {Write full present, imperfect, and future paradigms of \Latin{sum} and
  \Latin{possum}.  Then parse and translate \Latin{puer in horto est et
  laborare potest, sed hodie quiescit}.}
  {Explain what construction \Latin{potest} opens and what completes it.}
  {The paradigms should match the packet tables.  Translation: the boy is in
  the garden and can work, but today he rests.  \Latin{potest} is completed by
  the present active infinitive \Latin{laborare}.}
\SubmoduleExtension
  {Present-system contrast}
  {Present, imperfect, and future forms distinguish ongoing/current,
  continuing past, and future viewpoints.}
  {For \Latin{porto, moneo, mitto, audio}, form third-person singular and
  plural in all three present-system tenses.  Use six of the forms in a short
  account of work in a town.}
  {Without notes, mark conjugation and personal ending in four randomly chosen
  forms.}
  {Expected stems and tense signs follow the packet paradigms: for example
  \Latin{portat/portant, portabat/portabant, portabit/portabunt};
  \Latin{mittit/mittunt, mittebat/mittebant, mittet/mittent}.  The account is
  sound when time relations and subject number remain consistent.}
\SubmoduleExtension
  {A first sustained paragraph}
  {Repeated vocabulary frees attention for verb form, clause connection, and
  reference.}
  {Expand to six sentences: \Latin{Mane agricola surgit. Ad agrum ambulat.
  Ibi laborat.}  Include one imperfect background, one future plan, one
  sentence with \Latin{possum}, and two adjectives.}
  {Cover the draft and rebuild its finite-verb spine in order.}
  {Many paragraphs are possible.  A checkable model: \Latin{Mane agricola
  diligens surgit. Ad agrum magnum ambulat. Ibi cotidie laborabat. Hodie
  fossam parat. Cras semen portabit. Multum laborare potest.}  Each finite
  verb must agree with its subject and express the stated time relation.}
\end{courseunit}

\begin{courseunit}{M05}
\SubmoduleExtension
  {Completed action across the perfect system}
  {The perfect stem supports perfect, pluperfect, and future perfect forms;
  their temporal relations matter more than a one-word English gloss.}
  {Build all three active perfect-system tenses of \Latin{video} and
  \Latin{scribo}.  Use a pluperfect as background for a perfect event and a
  future perfect as the condition preceding a future action.}
  {Recover each form's stem, tense marker, and personal ending.}
  {Models include \Latin{vulpem videram antequam canis cucurrit} and
  \Latin{cum epistulam scripseris, eam mittes}.  Other propositions are sound
  when the perfect-system relations and morphology are accurate.}
\SubmoduleExtension
  {Passive form and deponent meaning}
  {Ordinary passive verbs change the subject's relation to the action;
  deponents retain passive-looking forms with active lexical meaning.}
  {Parse and translate \Latin{porta aperitur}, \Latin{porta aperta est},
  \Latin{viator ducem sequitur}, and \Latin{viator ducem secutus est}.  Turn
  the first two into active clauses with an expressed agent.}
  {State what morphology and lexicon each contribute to a deponent parse.}
  {Translations: the gate is opened/is being opened; the gate was opened; the
  traveler follows the guide; the traveler followed the guide.  Active models:
  \Latin{custos portam aperit; custos portam aperuit}.  Morphology identifies
  the passive-looking forms; the lexical entry identifies active deponent
  meaning.}
\SubmoduleExtension
  {Narrative tense and voice ledger}
  {A connected account must keep event order, participant reference, and voice
  distinct.}
  {Write an eight-line ordinary-world narrative containing present,
  imperfect, perfect, pluperfect, future, one ordinary passive, and one
  deponent.  Beside it, make a ledger of every finite form and its job.}
  {Retell the account from memory with the same events but a different Latin
  order.}
  {No single wording controls.  The model check asks whether every listed
  tense establishes the intended time relation, every passive preserves
  agent/patient roles, every deponent is translated actively, and all subjects
  agree with their verbs.}
\end{courseunit}

\begin{courseunit}{M06}
\SubmoduleExtension
  {Reference and government in motion}
  {Pronouns sustain reference across sentences; prepositions and verbs govern
  forms independently of English word order.}
  {Annotate every pronoun and governed case in \Latin{Marcus ad villam venit.
  Ibi amicum suum videt et cum eo ad flumen ambulat. Quod prope villam fluit,
  id latum est}.  Repair any ambiguous reference.}
  {Give the antecedent, case, number, gender, and function of each pronoun.}
  {\Latin{ibi} refers to the villa; \Latin{suum} refers to Marcus and agrees
  with \Latin{amicum}; \Latin{eo} is ablative after \Latin{cum}.  As written,
  \Latin{quod/id} most naturally refers to neuter \Latin{flumen}; a sound
  annotation notes that immediate agreement resolves the chain.}
\SubmoduleExtension
  {Seven-pass connected reading}
  {A repeatable workflow separates observation, morphology, syntax, literal
  structure, natural sense, and later lookup.}
  {Apply the packet's complete reading sequence to the ordinary paragraph in
  the reader.  Keep separate ledgers for printed forms, supplied words,
  uncertain lemmas, and lookup results.}
  {Close all references and summarize the paragraph directly from the Latin.}
  {The model is the reader's keyed parse.  A complete process never records a
  supplied English helper as though it were printed Latin and never changes an
  initial hypothesis without recording the evidence.}
\SubmoduleExtension
  {From ordinary prose toward the Missal}
  {The same morphology and clause map must work in neutral prose and in a short
  received liturgical passage.}
  {Compare one reader paragraph with the packet's Missal worked passage.  Mark
  the finite spine, governed forms, reference chains, and every compression or
  ellipsis in both.  Compose a neutral four-sentence paragraph using two of the
  same constructions but none of the liturgical wording.}
  {Explain which difficulty belongs to Latin generally and which belongs to
  the liturgical genre.}
  {The comparison should identify shared morphology and syntax before genre
  features.  The project paragraph must be labeled instructional and must not
  present itself as a prayer or received formula.}
\end{courseunit}

\begin{courseunit}{M07}
\SubmoduleExtension
  {Reference chains and comparison}
  {Relative, interrogative, and indefinite forms organize reference while
  comparative forms organize degree.}
  {Parse and translate \Latin{Viator qui celerius ambulat prius ad urbem
  venit; ille autem quem videmus lentissimus est}.  Identify each antecedent,
  comparison, and adverb.}
  {Replace \Latin{viator} with a feminine plural subject and make every
  dependent form agree.}
  {Translation: the traveler who walks more swiftly reaches the city earlier;
  but that man whom we see is slowest.  A feminine plural recast may begin
  \Latin{viatrices quae celerius ambulant...}; the remaining clause must be
  revised according to the intended referent.}
\SubmoduleExtension
  {Number as morphology and information}
  {Cardinals, ordinals, distributives, and numerical adverbs answer different
  questions and behave differently in syntax.}
  {Write six ordinary sentences distinguishing \English{three}, \English{the
  third}, \English{three each}, and \English{three times}; include one date or
  distance expression and parse every numeral.}
  {Retrieve the irregular comparison of \Latin{bonus}, \Latin{malus},
  \Latin{magnus}, and \Latin{parvus}.}
  {Models will vary.  In a sound model the numeral class matches the
  proposition and agrees or governs as taught.  Irregular series include
  \Latin{bonus, melior, optimus}; \Latin{malus, peior, pessimus};
  \Latin{magnus, maior, maximus}; \Latin{parvus, minor, minimus}.}
\SubmoduleExtension
  {Expository composition seed}
  {Definition and comparison are early building blocks of later argumentative
  and periodic prose.}
  {Define two kinds of road, compare them by three features, and conclude
  which a traveler chooses under one stated condition.  Use a relative clause,
  one comparative adjective, one comparative adverb, and an ordinal.}
  {Underline the exact Latin that connects each claim.}
  {A sound paragraph states real propositions, gives each relative an explicit
  antecedent, completes comparisons when needed, and uses the ordinal as an
  ordering expression rather than as a cardinal.}
\end{courseunit}

\begin{courseunit}{M08}
\SubmoduleExtension
  {Infinitive time and voice}
  {In indirect statement an accusative subject and infinitive express a
  reported proposition; infinitive tense is normally relative to the
  governing verb.}
  {Turn \Latin{nauta venit}, \Latin{nauta veniebat}, and \Latin{nauta veniet}
  into reports after \Latin{dicit} and after \Latin{dixit}.  Explain every
  change.}
  {From memory, form present, perfect, and future active infinitives of one
  verb and the available passive counterparts.}
  {Models after \Latin{dicit}: \Latin{nautam venire}, \Latin{nautam venisse},
  and \Latin{nautam venturum esse}; the same infinitives after \Latin{dixit}
  retain relative simultaneity, anteriority, and posteriority.  Context may
  alter the best English tense.}
\SubmoduleExtension
  {Participles attached and absolute}
  {A participle agrees with its controller; an ablative absolute forms a
  grammatically detached circumstance.}
  {Parse \Latin{sole oriente agricolae laborant; vulpes a cane visa fugit;
  viator epistulam scriptam legit}.  Distinguish every participle's tense,
  voice, controller, and clause relation.}
  {Recast one attached participle as a finite subordinate clause.}
  {\Latin{oriente} agrees with ablative \Latin{sole} in an absolute;
  \Latin{visa} is perfect passive and agrees with subject \Latin{vulpes};
  \Latin{scriptam} is perfect passive and agrees with object \Latin{epistulam}.
  Finite recasts must preserve relative time and voice.}
\SubmoduleExtension
  {Narrative compression and expansion}
  {Non-finite forms can compress propositions, but an accurate reader can
  expand them without inventing agents or time relations.}
  {Compress a six-clause project narrative using one indirect statement, two
  different participles, and an ablative absolute.  Then expand the result
  back into finite clauses and compare propositions.}
  {Circle any information introduced only during expansion.}
  {A defensible pair preserves the same agents, patients, time relations, and
  negation.  Any supplied conjunction or agent belongs in an explicit ledger;
  silent additions show where the compression was not actually understood.}
\end{courseunit}

\begin{courseunit}{M09}
\SubmoduleExtension
  {The complete subjunctive system}
  {Recognition begins with morphology; clause classification follows from the
  governor, connective, sequence, and discourse force.}
  {Form and parse all four subjunctive tenses of \Latin{laudo}, \Latin{moneo},
  \Latin{mitto}, and \Latin{audio} in one person assigned by the memory sheet.
  Add the corresponding forms of \Latin{sum}.}
  {State primary and historic sequence without using an English translation as
  proof.}
  {Forms should match the packet paradigms.  Primary sequence normally pairs a
  primary governor with present or perfect subjunctive; historic sequence
  normally pairs a secondary governor with imperfect or pluperfect.  Relative
  time and construction still govern interpretation.}
\SubmoduleExtension
  {Purpose, result, command, and substantive work}
  {The same subjunctive form can serve different constructions; signals such
  as \Latin{tam}, \Latin{ut non}, a verb of commanding, or a noun-like clause
  help decide.}
  {Classify and translate four packet examples, then compose one new example
  each of purpose, result, indirect command, and substantive clause.  Label
  governor, signal, negation, sequence, and clause function.}
  {Convert purpose to negative purpose and result to negative result.}
  {Negative purpose ordinarily uses \Latin{ne}; negative result uses
  \Latin{ut non}.  Each composition must contain an appropriate governor and
  express a proposition suited to its named construction, not merely include
  \Latin{ut} plus a subjunctive.}
\SubmoduleExtension
  {Composition through subordination}
  {Later prose depends on deliberately relating propositions rather than
  translating English connective by connective.}
  {Plan five propositions about building or repairing a house.  Join them in
  a paragraph using purpose, result, indirect command, and one substantive
  clause.  Then rewrite the paragraph with a historic governing verb and make
  the necessary sequence changes.}
  {Read only the finite verbs and conjunctions to test the architecture.}
  {A sound pair preserves the propositions while revising dependent tenses
  where sequence calls for it.  Every subordinate clause has a recoverable
  governor and function; mere tense substitution without a coherent time
  relation is not a model.}
\end{courseunit}

\begin{courseunit}{M10}
\SubmoduleExtension
  {Circumstance from evidence}
  {Temporal, causal, and concessive clauses are classified from connective,
  mood, context, and the relation of their propositions.}
  {For six packet clauses, make a grid of connective, mood, tense, proposition,
  and relation.  Then compose a matched set in which one event is presented as
  time, cause, and concession.}
  {Explain why the connective alone cannot always settle the analysis.}
  {The matched set should genuinely change the relation asserted, for example
  \Latin{cum pluit...} as time or circumstance, an explicit causal ground, and
  a concession whose main clause remains true despite it.  Mood and context
  must support the label.}
\SubmoduleExtension
  {Conditions and counterfactual thought}
  {Protasis and apodosis form one logical pair; tense and mood present the
  condition as open, vivid, potential, or contrary to fact.}
  {Write present and past contrary-to-fact conditions from the same
  proposition, then open and future-more-vivid versions.  Parse every finite
  verb and state what the writer commits to.}
  {Change one condition's polarity without losing its logical relation.}
  {Expected contrasts use indicative forms for open conditions as appropriate,
  present subjunctive for potential force, imperfect subjunctive for present
  contrary to fact, and pluperfect subjunctive for past contrary to fact.
  Models must be judged by the whole pair, not one verb alone.}
\SubmoduleExtension
  {Characteristic and factual relative clauses}
  {A relative clause may identify a fact about a definite antecedent or
  characterize a kind of person or thing.}
  {Contrast \Latin{hominem video qui laborat} with \Latin{nemo est qui semper
  laboret}.  Create four further pairs that keep vocabulary stable while
  changing factual identification into characteristic description.}
  {Name antecedent, mood, and discourse claim in each pair.}
  {The indicative example reports work by a particular man; the subjunctive
  characterizes the non-existent or generalized class.  New pairs need an
  antecedent and context that actually support the distinction; changing mood
  alone is insufficient.}
\end{courseunit}

\begin{courseunit}{M11}
\SubmoduleExtension
  {Commands, prohibitions, wishes, and fear}
  {Independent and dependent volitional language uses several constructions;
  person, negation, and governor distinguish them.}
  {Transform one project proposition into a direct command, prohibition,
  indirect command, present wish, past unfulfilled wish, and fear clause.
  Label construction and decisive form.}
  {Retrieve the negation ordinarily used in each construction.}
  {Models will vary.  Commands may use the imperative or jussive subjunctive;
  prohibitions may use \Latin{noli/nolite} plus infinitive or \Latin{ne} with
  the appropriate subjunctive; wishes use \Latin{utinam}; fear clauses reverse
  the easy English expectation, with \Latin{ne} fearing that and \Latin{ut}
  fearing that not.}
\SubmoduleExtension
  {Gerund, gerundive, supine, and obligation}
  {Verbal nouns and adjectives preserve verbal government while taking nominal
  or adjectival form; the passive periphrastic expresses necessity with a
  dative of agent.}
  {Parse and translate \Latin{ars scribendi}, \Latin{ad urbem videndam},
  \Latin{mirabile dictu}, and \Latin{nobis liber legendus est}.  Transform a
  gerund with object into a gerundive construction.}
  {Explain what agrees in each construction and what retains verbal force.}
  {Translations include the art of writing; for seeing/visiting the city;
  marvelous to say; the book must be read by us.  A gerundive attraction makes
  the verbal adjective agree with its former object while retaining the
  intended relation.}
\SubmoduleExtension
  {Instructional prose and obligation}
  {Commands and periphrastics prepare a reader for rubrical, legal, and
  scholastic prose as well as for ordinary instructions.}
  {Write eight connected instructions for preparing a journey.  Use two direct
  commands, one prohibition, one indirect command, a gerundive of purpose, a
  passive periphrastic, and one wish.  Then recast the passage as reported
  instructions.}
  {Make a government ledger for every verbal noun or adjective.}
  {A checkable response preserves persons and propositions during the recast,
  uses each construction for its proper discourse job, and distinguishes
  obligation from a simple future prediction.}
\end{courseunit}

\begin{courseunit}{M12}
\SubmoduleExtension
  {Cases used adverbially}
  {Place, time, means, manner, and cause may be expressed by case alone or by a
  prepositional phrase; the choice is lexical and idiomatic, not mechanical.}
  {Sort twenty packet phrases by form and function.  Then write an ordinary
  travel paragraph containing place where, place to which, place from which,
  duration, time when, means, manner, and cause.}
  {For each phrase, state why a preposition appears or does not appear.}
  {The packet tables and key control the sorts.  In the paragraph every phrase
  must be parsed before its English label is assigned; a plausible translation
  alone does not establish the Latin construction.}
\SubmoduleExtension
  {Order, ellipsis, attraction, and compression}
  {A marked order may foreground material; ellipsis omits recoverable words;
  attraction aligns a form with nearby structure; compression combines these
  without suspending grammar.}
  {Take a short packet passage and produce (1) a diplomatic word list, (2) a
  neutral-order scaffold, (3) a supply ledger, and (4) a natural translation.
  Mark every change of order and every supplied word.}
  {Restore the original Latin from the scaffold without looking.}
  {A sound scaffold preserves every printed form and proposition.  Supplied
  words remain bracketed and justified; the neutral order is an analytic aid,
  never a claim that the received author wrote it.}
\SubmoduleExtension
  {Core-grammar synthesis}
  {Reading unfamiliar prose calls for the entire nominal, verbal, subordinate,
  and discourse system at once.}
  {Analyze the packet's authentic prose without translating for the first
  pass.  Then write a structural translation, a natural translation, and a
  six-sentence imitation that changes every proposition while preserving two
  named syntactic frames.}
  {List which forms or constructions still need deliberate lookup.}
  {The keyed reader supplies the parse and translations.  The imitation must
  be visibly project-composed, preserve the declared frames rather than copied
  wording, and withstand a morphology/government/reference check.}
\end{courseunit}

\begin{courseunit}{M13}
\SubmoduleExtension
  {Independent workflow and the fixed Ordinary}
  {A familiar liturgical sequence can assist navigation but must not replace
  grammatical reading.}
  {On a fresh Ordinary passage, record text class, finite spine, clause bounds,
  morphology, government, supplied words, literal structure, natural sense,
  and only then lookup results.  Repeat the process on the paired neutral or
  classical paragraph.}
  {State one place where familiarity tempted you to skip evidence.}
  {The packet keys control the local parse.  A complete record distinguishes
  what the page prints from memory, rubric, translation, and editorial supply;
  it applies the same grammatical standard to both genres.}
\SubmoduleExtension
  {Formula and free syntax}
  {Fixed formulas contain ordinary Latin syntax even when repetition makes
  their meaning immediately familiar.}
  {Choose four Ordinary formulas.  Expand each into a full clause where that is
  defensible, mark what is supplied, and compose an ordinary-world sentence
  using the same grammatical frame but unrelated vocabulary.}
  {Close the Missal text and parse the four formulas from form alone.}
  {Models accept more than one supply when the printed phrase permits it, but
  every addition is bracketed.  The neutral sentences test the construction
  without presenting project wording as liturgical.}
\SubmoduleExtension
  {Reading memorandum}
  {Direct comprehension is strengthened by concise Latin summary and by an
  explicit uncertainty ledger.}
  {Write a five-sentence Latin memorandum of the two passages: two propositions
  from each and one precise comparison of style.  Record every borrowed idiom
  and every unresolved question.}
  {Revise the memorandum once for clause attachment and once for morphology.}
  {A sound memorandum reports the source propositions accurately, separates
  grammatical observation from theological inference, and documents rather
  than conceals borrowed wording or uncertainty.}
\end{courseunit}

\begin{courseunit}{M14}
\SubmoduleExtension
  {The architecture of petition}
  {Collects, Secrets, and Postcommunions compress address, ground, petition,
  offering or sacramental context, and desired effect into tightly governed
  periods.}
  {Map three packet orations by finite spine, address, relative or participial
  ground, petition, purpose/effect, and conclusion.  Leave a label open where
  grammar permits more than one discourse relation.}
  {Rebuild each sentence from its clause map without the English.}
  {The packet key controls the forms.  A defensible map follows governors and
  clause attachment; it does not force every \Latin{ut} clause into the same
  semantic category or treat a conventional conclusion as the main petition.}
\SubmoduleExtension
  {Petition and periodic prose outside the Missal}
  {Cicero and Augustine offer longer periodic organization against which the
  distinctive compression of an oration can be seen.}
  {Compare one packet oration with the paired authentic prose excerpt.  Mark
  postponed completions, parallel members, connective chains, and the location
  of the main assertion.}
  {Mark each period's visual phrasing, including its main assertion, dependent
  members, and delayed completions.}
  {The comparison should identify observable syntax and rhetoric, not merely
  call one text ornate.  Main assertions, dependent members, and parallels
  must be cited from the printed Latin.}
\SubmoduleExtension
  {Composition from propositions, not formula theft}
  {Earlier sentence transformations now support a short, coherent, explicitly
  instructional composition.}
  {Plan an address-free civic petition about repairing a road: state condition,
  request, ground, and intended effect.  Draft it first as four sentences and
  then as one period.  Document any borrowed construction.}
  {Compare the one-period version with the four-sentence version for lost or
  distorted propositions.}
  {A model keeps the civic, project-composed identity visible; preserves every
  proposition; gives each subordinate clause a governor; and records any
  source-supported frame without copying a liturgical formula wholesale.}
\end{courseunit}

\begin{courseunit}{M15}
\SubmoduleExtension
  {Argument and narrative on the page}
  {Lessons and Epistles develop claims and inferences; Gospels organize events,
  participants, speech, and biblical coordination.}
  {Make an argument map for the packet Epistle and an event/speech map for the
  Gospel.  For both, list finite clauses, connectives, pronoun antecedents, and
  shifts of speaker or time.}
  {Summarize each passage directly in three Latin sentences.}
  {The keyed analysis controls local details.  The argument map must distinguish
  support, contrast, inference, and conclusion; the narrative map must retain
  event order and speakers without importing details from memory.}
\SubmoduleExtension
  {Cicero, Augustine, and biblical prose in contrast}
  {Different prose traditions may use the same grammar with different
  connective density, balance, repetition, and word order.}
  {Annotate the paired authentic excerpts for coordination, subordination,
  periodic delay, repetition, and explicit argument markers.  Cite one feature
  each that belongs to an author or passage, not to Latin universally.}
  {Translate one sentence structurally and then naturally, preserving a
  documented ambiguity.}
  {A sound comparison cites exact forms and does not infer chronology,
  theology, or influence from style alone.  The translations show the same
  proposition and identify any material alternative in a note.}
\SubmoduleExtension
  {Narrative and argument composition}
  {Composition tests whether syntax can be selected for a purpose rather than
  merely recognized.}
  {Write one six-sentence narrative and one six-sentence argument about the
  same ordinary event.  Reuse the core vocabulary but change clause relations,
  connectives, tense structure, and paragraph architecture.}
  {Create a source/idiom ledger for both drafts and complete three revision
  passes: structure, concord, idiom.}
  {The two pieces must differ by discourse work, not only by vocabulary.
  Every claim or event remains recoverable; reference is stable; and any
  source-shaped phrase is documented.}
\end{courseunit}

\begin{courseunit}{M16}
\SubmoduleExtension
  {Poetic order in the main course}
  {Chants and sequences use parallelism, ellipsis, imagery, and marked order;
  a reader can recover their syntax before undertaking formal prosody.}
  {For a packet chant or sequence, mark clause boundaries, parallel members,
  agreeing pairs separated by order, omitted recoverable words, and images.
  Produce a neutral-order scaffold without altering the received text.}
  {Name the evidence joining every separated pair.}
  {The packet key controls the parse.  The scaffold is a project annotation,
  not a rewritten source; it must preserve every printed form and record each
  supply.  Formal scansion belongs to the later poetry branch.}
\SubmoduleExtension
  {The long prose period}
  {Prefaces and the Canon sustain grammatical dependencies across balanced and
  formulaic members.}
  {Bracket one long period into main assertion, relative or participial
  expansions, coordinated members, and conclusion.  Track every relative and
  pronoun to its controller.}
  {Reconstruct the finite spine with all modifiers removed, then replace them.}
  {A complete bracket map identifies the governing finite form and attaches
  every dependent member once.  Removing material may yield inelegant Latin;
  the exercise is analytical and may not change the received transcription.}
\SubmoduleExtension
  {Prose--poetry transfer}
  {Reading competence includes moving between dense prose and verse without
  assuming that verse has no syntax.}
  {Choose one proposition from the Preface and one from the poem.  Restate each
  in plain project prose, then compare what order, parallelism, written pattern
  visible on the page, and imagery contribute.}
  {Write four lines of rhythmic project Latin only as syntactic imitation; do
  not claim meter before the advanced branch.}
  {The prose restatements preserve propositions and remain labeled project
  compositions.  The comparison rests on visible textual features; it neither
  supplies a pronunciation nor calls unscanned lines metrical.}
\end{courseunit}

\begin{courseunit}{M17}
\SubmoduleExtension
  {Reading across the Temporal}
  {Seasonal familiarity helps locate texts but cannot determine a form's
  morphology or a passage's exact source.}
  {Read two unfamiliar Temporal selections using only their printed headings
  and Latin.  Record genre, text class, finite spine, governing construction,
  and what lookup changed.}
  {Explain one rejected guess and the evidence against it.}
  {The course inventory controls source identity.  A complete record separates
  heading evidence, grammar, remembered association, and verified lookup; a
  correct seasonal guess without the Latin analysis is not sufficient evidence.}
\SubmoduleExtension
  {Sanctoral and Common without formula substitution}
  {Repeated Common language creates useful lexical recurrence but also tempts
  the reader to substitute a remembered formula for the printed one.}
  {Collate a Sanctoral passage with its related Common.  Mark exact identity,
  inflectional change, lexical substitution, omission, and addition.  Parse
  both witnesses before interpreting the difference.}
  {Cover one witness and reconstruct only what is actually shared.}
  {A sound collation preserves both bases, records each difference once, and
  never silently corrects one witness from the other.  Interpretation follows
  rather than replaces the textual record.}
\SubmoduleExtension
  {Register comparison memorandum}
  {The Missal's seasonal and sanctoral registers can be described more
  accurately when compared with non-liturgical prose already read.}
  {Write a one-page memorandum comparing one Missal passage, one classical
  passage, and one patristic passage.  Discuss clause length, connective use,
  order, lexical recurrence, and compression with exact examples.}
  {Delete every adjective of praise or dislike that is not tied to evidence.}
  {A defensible memorandum identifies text and locus, cites observable Latin,
  distinguishes authorial from genre claims, and leaves broader historical
  explanations open unless separately sourced.}
\end{courseunit}

\begin{courseunit}{M18}
\SubmoduleExtension
  {Holy Week and Requiem registers}
  {Distinctive ceremonies and remembered translations must not displace the
  local Latin text, genre, and grammatical evidence.}
  {Analyze one Holy Week and one Requiem passage through text class, finite
  spine, compression, biblical texture, and natural sense.  Record what cannot
  be decided by grammar alone.}
  {Summarize each directly from the Latin before consulting a translation.}
  {The packet's answer key, reader apparatus, and source notes control local
  decisions.  A complete response does not infer ceremonial directions,
  doctrine, or textual history from morphology and labels such questions for
  their proper sources.}
\SubmoduleExtension
  {Complete-Mass sustained reading}
  {The destination skill is continuous, evidence-led reading across changing
  Missal genres, supported by the broad Latin competence built beforehand.}
  {Work through one of the complete formularies printed after the M18 reader
  across as many returns as are useful. Maintain
  a running vocabulary ledger, clause map, source-class record, uncertainty
  list, structural translation, and concise Latin summaries of each genre.}
  {Return after an interval and reread selected pages without the English or
  prior annotations.}
  {Before continuing, one should reasonably be able to sustain the process,
  recover most grammar
  from forms, recognize when lookup is needed, and correct recurring errors
  from the supplied apparatus.}
\SubmoduleExtension
  {Beyond one familiar corpus}
  {A reader of the Missal should also be able to approach other Latin without
  depending on memorized liturgical wording.}
  {Read the Cicero definition and the project-composed bridge narrative printed
  in the M18 reader.  Use the same workflow, then write a comparison of what
  transfers and what is register-specific.}
  {Choose the next path: broader prose and composition in M19, or continued
  consolidation from the memory workbook before doing so.}
  {A sound response uses morphology and syntax consistently across all three
  texts, documents unresolved lexical or historical questions, and describes
  register from cited forms rather than from reputation.}
\end{courseunit}

\begin{courseunit}{M19}
\SubmoduleExtension
  {Register across time}
  {Classical, biblical, patristic, scholastic, and liturgical Latin overlap but
  do not form a single undifferentiated style.}
  {Build a register dossier from the packet's named passages: exact locus,
  clause architecture, coordination/subordination, marked order, recurring
  lexicon, and one negative result.  Avoid assigning dates or causes from a
  feature alone.}
  {Rewrite one sentence in neutral project prose and list every change.}
  {A defensible dossier distinguishes observation from historical explanation,
  records absent as well as present expected features, and labels the rewrite
  as project-composed rather than as an improvement of the source.}
\SubmoduleExtension
  {Biblical texture and source-language restraint}
  {Repetition, coordination, and certain constructions can describe biblical
  texture; a claim of source-language influence needs direct comparative
  evidence.}
  {Use the Matthew 8:8 Vulgate/Ordinary pair printed in A2 and set it beside
  the Ambrose passage in the packet reader.  Mark observable texture, exact
  adaptation, and merely possible parallel.  State which claims the evidence
  does not establish.}
  {Produce literal and natural translations that preserve one material
  ambiguity.}
  {The model keeps quotation identity, grammatical description, and historical
  explanation in separate columns.  It does not call every coordination a
  Hebraism or infer dependence from a shared commonplace.}
\SubmoduleExtension
  {Composition in two registers}
  {Advanced composition begins by controlling propositions and syntax before
  making restrained stylistic choices.}
  {Write one 150-word ordinary exposition, then recast it in a consciously
  different but defensible Latin register using two named source features.
  Keep a sentence-by-sentence proposition and idiom ledger.}
  {Reverse two stylistic changes and test whether the propositions remain.}
  {A sound pair is grammatically complete, documents every source-shaped idiom,
  changes style rather than factual content, and avoids pastiche made from
  copied phrases.}
\end{courseunit}

\begin{courseunit}{M20}
\SubmoduleExtension
  {Allusion, quotation, and parallel}
  {Exact quotation, adaptation, formula, echo, and thematic parallel make
  different evidence claims.}
  {For four packet pairs, transcribe both loci, align exact words and forms,
  classify the relation, record differences, and give a confidence statement.
  Include one honest negative result.}
  {Remove all source labels and see whether the verbal evidence alone supports
  the classification.}
  {The model demands exact loci and wording.  Shared theme without distinctive
  verbal contact remains a parallel; resemblance may be inconclusive; a
  negative result is preferable to an invented source claim.}
\SubmoduleExtension
  {Ciceronian period and Augustinian sentence sequence}
  {Periodic prose suspends or balances completion through syntax; a linked
  sequence can instead build force across complete sentence boundaries.
  Rhetoric is the work performed by those choices, not a list of figure names.}
  {Diagram Cicero's complete definition sentence and Augustine's complete
  three-sentence opening: main assertions, sentence, colon, and member
  boundaries, anaphora or balance, delayed elements, climax, cross-sentence
  repetition, and rhythm visible as textual arrangement without pronunciation
  claims.}
  {Turn Cicero's sentence and each Augustine sentence into short neutral
  clauses, then state what is lost.}
  {A defensible diagram attaches every member grammatically, preserves the
  three Augustine sentence boundaries, and ties each rhetorical observation
  to an exact form or order.  The neutral restatement may lose emphasis,
  suspense, balance, linkage, or cadence while preserving propositions.}
\SubmoduleExtension
  {Periodic imitation}
  {Imitation becomes composition when new propositions are organized by a
  declared structure rather than filled into copied wording.}
  {Extract an abstract frame from one source period.  Plan new civic or moral
  propositions, draft a 120--180 word period or paired periods, and annotate
  every point of structural imitation and departure.}
  {Rewrite the draft as short sentences to audit logical completeness.}
  {The model check asks whether the abstract frame is accurately named,
  every proposition is original and recoverable, morphology and government
  are sound, and source wording has not been silently appropriated.}
\end{courseunit}

\begin{courseunit}{M21}
\SubmoduleExtension
  {Diplomatic text and normalization}
  {Textual comparison begins by preserving each witness before regularizing
  orthography or punctuation.}
  {Transcribe the supplied base representation diplomatically, create a
  separately labeled normalized copy, and log every changed character,
  abbreviation, punctuation mark, capitalization, and line boundary.}
  {Reproduce three lines from the diplomatic copy without looking at the
  normalized text.}
  {The reference transcription controls the check.  The two layers remain
  separate; uncertain characters are marked rather than guessed; normalization
  never silently changes morphology, word division, or source evidence.}
\SubmoduleExtension
  {Collation and commentary}
  {A collation records witness differences; commentary may then discuss their
  grammatical or interpretive consequences without constructing an edition it
  cannot support.}
  {Collate two packet witnesses, classify every difference, and write notes on
  only the variants that affect form, syntax, sense, or attribution.  State the
  base and what the exercise does not establish.}
  {Reverse the base and verify that the same differences remain visible.}
  {A sound apparatus is complete and directionally consistent, distinguishes
  graphic from substantive variation, and does not call one reading original,
  corrupt, or superior without appropriate textual evidence.}
\SubmoduleExtension
  {Controlled imitation with an evidence ledger}
  {Imitation joins close reading to composition when the imitated feature can
  be named, reproduced with new content, and checked.}
  {Choose one clause frame, one ordering device, and one rhetorical pattern
  from Cicero, Augustine, or another named packet source.  Compose three new
  sentences for each and annotate both correspondence and divergence.}
  {Remove the source and explain the feature from your own examples.}
  {The examples must be independent propositions with accurate Latin.  The
  ledger cites exact source loci, avoids copying distinctive phrases, and
  treats failure to preserve a feature as information for revision.}
\end{courseunit}

\begin{courseunit}{M22}
\SubmoduleExtension
  {Guided composition from a proposition plan}
  {A sustained draft grows from explicit propositions, relations, genre, and
  source choices rather than from English sentence-by-sentence substitution.}
  {Plan and write 250--400 Latin words in one chosen prose form: narrative,
  exposition, commentary, civic petition, or source-based meditation clearly
  labeled as instructional.  Preserve the proposition plan and first draft.}
  {Map every finite clause back to the plan and identify additions or losses.}
  {No single text controls.  A mature draft has a recoverable architecture,
  stable reference, accurate morphology and government, coherent vocabulary,
  and a source/idiom ledger for every non-obvious borrowed frame.}
\SubmoduleExtension
  {Three-pass revision and counter-reading}
  {Separate reviews of clause structure, concord/reference, and idiom/style
  reveal different error families.}
  {Mark the first draft in three passes.  Then write a hostile paraphrase of
  each sentence's proposition to expose ambiguity, repair what the Latin did
  not actually say, and preserve a change log.}
  {After an interval, reconstruct one paragraph from its proposition plan.}
  {A sound revision explains each substantive change by rule or evidence,
  distinguishes correction from preference, and does not erase the earlier
  draft.  Reconstruction tests productive control rather than visual memory.}
\SubmoduleExtension
  {Independent research and the poetry branch}
  {Advanced prose competence supports responsible source work; formal verse
  asks for a distinct body of prosodic knowledge and therefore branches here.}
  {Complete one bounded research note on a text in the supplied anthology,
  recording its exact source, edition, locus, transcription class,
  claim/evidence separation, translation, and unresolved questions.  Then
  review the stated abilities at the opening of P1 before deciding whether to
  begin formal poetry.}
  {Write a concise account of what can and cannot be concluded from grammar
  alone.}
  {The research note should make the evidence locatable, distinguish source
  from project work, expose the reasoning, and preserve remaining uncertainty.
  P1 is a branch from this common trunk, not a
  condition imposed on further prose reading or composition.}
\end{courseunit}

  \endgroup}
